% =============================================================================
%  扉页、摘要、译者说明
% =============================================================================
\begin{titlepage}
\centering
\vspace*{3.0cm}
{\Huge\bfseries 为物理学家讲述的\\[0.25em] 朗兰兹猜想\par}
\vspace{1.0em}
{\large\itshape Langlands' Conjectures for Physicists\par}
\vspace{2.4em}
{\large Mark Goresky\par}
\vspace{0.4em}
{\normalsize 高等研究院（I.A.S.）讲座，2004 年 3 月 8 日；会后扩写稿\par}
\vspace{3.0em}
{\large 中译注本\par}
\vspace{0.6em}
{\normalsize 数译 · Langlands 导读系列 · SY02-42\par}
\vfill
{\small
\begin{tabular}{@{}rl@{}}
底本 & Mark Goresky, \emph{Langlands' Conjectures for Physicists},\\
     & I.A.S. 讲座（2004-03-08）扩写稿，20 页（dvips 排版）\\
译者 & 数译（DSH agent：lead）\\
版本 & v1.0，2026 年 11 月 27 日\\
体例 & 公式编号 (2.0.1)、(4.2.1) 等照原文；\\
     & 原文 28 条尾注照录为文末「注释」；\\
     & 译注（橙色框）与「本节重点」（蓝色框）为译者所加\\
\end{tabular}}
\end{titlepage}

\section*{摘要}
\addcontentsline{toc}{section}{摘要}

这份讲义是作者 2004 年在普林斯顿高等研究院给一群物理学家讲的几次课的扩写稿。
它要做的事情只有一件：把朗兰兹猜想这句话

\begin{center}
\itshape
「$\Gal(\Qbar/\bbQ)$ 的表示」对应「$\GL(n,\bbA_{\bbQ})$ 的自守表示」，
\end{center}

\noindent
从零拆开，逐词解释，然后再把每一块拼回去。

全篇分八节。第 1 至 3 节讲背景：域的种类（数域、$p$ 进域、函数域），以及「在某个域上解方程」这件事在不同域上意味着什么。
第 4 节讲左边：$\ell$ 进平展上同调、Frobenius 元素、$L$ 函数。
第 5 节讲右边的最简单情形 $\sL(2)$：模形式、Hecke 算子、Hecke 代数。
第 6 节把右边换成现代语言：平代环、自守表示、Satake 参数。
第 7 节把两边对上：$n=2$ 时 Eichler–Shimura 已给出证明，$n\ge 3$ 时同一条路走不通——这是全篇最诚实的一句话。
第 8 节给出对任意群、任意域的推广形式，然后停在这里。

作者在开头就说明，这是一份尚未完成的工作：第 7.3 节与第 8.2–8.7 节只有标题，正文是空的。
本次翻译照实呈现这些空缺，没有代拟内容。

\section*{译者说明}
\addcontentsline{toc}{section}{译者说明}

\subsection*{一、这份译本的定位}

底本是 Mark Goresky 2004 年 3 月 8 日在 I.A.S. 为物理学家所作讲座的扩写稿，20 页，作者自印，未正式发表。
它在朗兰兹纲领的文献里不算前沿，也不算权威——作者自己说，它「充其量只是许多优秀综述的导引」。
那为什么还要译？

因为它恰好补上了本库最缺的一块：\textbf{从零到一的路线图}。
本库已有的几何朗兰兹材料（Gaitsgory 的证明大纲、Scholze 的 Bourbaki 报告、Sandvik 的 AB 等价讲义）
都假定读者已经知道「Frobenius 特征值对应 Hecke 特征值」这句话是什么意思。
这份讲义花了一半篇幅，就为了把这句话讲清楚：什么叫一个域的 Frobenius，什么叫一个模形式的 Hecke 特征值，
为什么它们会出现在同一个 $L$ 函数的两边，以及为什么 $n=2$ 能证、$n\ge 3$ 不能证。

更值得注意的是它的结尾。第 8.5 节标题是「几何 Langlands」，正文是空的。
这不是敷衍——2004 年的几何朗兰兹还没有今天这套话语，作者写到这里就停笔了。
于是这份讲义在结构上成了一座桥：桥这头是数论朗兰兹的完整入门，桥那头正好是本库工作的起点。
读它，等于把本库的选题动机重新走一遍。

\subsection*{二、体例}

\begin{enumerate}[leftmargin=2.2em,itemsep=2pt]
\item \textbf{编号}：公式编号一律沿用原文的「(节.小节.序号)」格式，如 (2.0.1)、(6.3.2)，
      用 \verb|\tag| 手工标注，不做重排，便于与原文左右对读。
\item \textbf{尾注}：原文把技术细节放在文末的 Notes 里，共 28 条。译文照录为末节「注释」，
      正文以上标数字标记；不改成脚注，以免破坏原文「正文浅、注释深」的两层结构。
\item \textbf{术语}：首次出现写作「中文（English）」，此后只用中文；
      与 \texttt{工程文件/work/glossary.tsv} 一致。数论与表示论的标准术语沿用国内通行译法
      （平展上同调、自守表示、Hecke 算子、Satake 参数、对偶群），
      没有通行译法的保留英文。
\item \textbf{译注}：只写三类内容——术语背景、原文含糊或省略之处、原文笔误。
      每条 1–5 句，不代替读者作证明，也不补写作者没写的内容。
\item \textbf{照实呈现}：原文自认未完成的地方（第 7.3 节、第 8.2–8.7 节）保留标题并标注「原文此处空缺」，
      不代拟、不删节。
\end{enumerate}

\subsection*{三、原文疑误一览}

下列各处按原文照录，并在正文相应位置以译注标出；这里集中列出，便于读者一眼看到「哪些地方不能照抄」。
（底本为 dvips 排版的自印稿，以下多为排印错误，不影响数学内容。）

\begin{itemize}[leftmargin=2.2em,itemsep=3pt]
\item \textbf{p.5}：a good cohomology theory 印作 an good。
\item \textbf{p.4}：polynomials 印作 polynomimals。
\item \textbf{p.9}：representation 印作 repesentation。
\item \textbf{p.11}：structure 印作 strucutre。
\item \textbf{p.11}：adeles 印作 elemenets（同一段）。
\item \textbf{p.12}：satisfies 印作 satisies。
\item \textbf{p.7}：(4.4.1) 式右侧印作 $\prod_p\zeta_p(p^{-s})$，漏了表示 $\rho$；
      按上下文（及前一式 $\zeta_p(\rho,t)$ 的定义）应为 $\prod_p\zeta_p(\rho,p^{-s})$。
\item \textbf{p.15}：interpreted 印作 interpereted。
\item \textbf{p.18}：注释 25 中 locally 印作 localy。
\item \textbf{p.12}：(6.3.2) 的分子印作 $\GL(n,\bbA)$，而 (6.3.1) 与 (6.2.2) 的分子都是 $\GL(n,\bbQ)$；
      按上下文 (6.3.2) 的分子应为 $\GL(n,\bbQ)$，此处系笔误。
\end{itemize}

\subsection*{四、配套材料}

\begin{itemize}[leftmargin=2.2em,itemsep=2pt]
\item \textbf{资料卡}：\texttt{datasets/goresky04/}（JSON，概念卡 + 推导卡，附渲染版 Markdown）。
\item \textbf{资料卡五件套}：\texttt{资料卡/00-速览.md} … \texttt{04-术语文献勘误.md}。
\item \textbf{术语表}：\texttt{工程文件/work/glossary.tsv}。
\end{itemize}

\subsection*{五、读法建议}

如果只为读懂本库其他材料，建议这样读：

\begin{enumerate}[leftmargin=2.2em,itemsep=2pt]
\item 先读第 2 节与第 6.5 节，一前一后，加起来不到一页，就是整个猜想的两句话。
\item 再读第 5 节。这一节是全篇最具体的地方：模形式、Hecke 算子、Hecke 代数、Fourier 系数，
      每一步都有公式可算。读完应能回答「一个模形式的 Hecke 特征值到底是什么」。
\item 然后读第 7.1 节（$n=2$ 为什么能证）与第 7.2 节（$n\ge 3$ 为什么卡住）。
      这两节加起来一页半，是理解几何朗兰兹动机的最短路径。
\item 第 3 节讲各种域，可以当作查表用；第 4 节讲平展上同调，物理读者第一遍可以跳过。
\item 尾注是全文的技术底层，注释 6（自守因子）与注释 25（Satake 同构）值得单独读。
\end{enumerate}
